\newcommand{\N}{\mathbb{N}}
\newcommand{\Z}{\mathbb{Z}}
\newcommand{\R}{\mathbb{R}}
\newcommand{\Q}{\mathbb{Q}}
\newcommand{\x}{\times}						
\newcommand{\id}{\mathrm{id}}
\newcommand{\Acal}{\mathcal{A}}
\newcommand{\Bcal}{\mathcal{B}}
\newcommand{\Ccal}{\mathcal{C}}
\newcommand{\Dcal}{\mathcal{D}}
\newcommand{\Ecal}{\mathcal{E}}
\newcommand{\Fcal}{\mathcal{F}}
\newcommand{\Gcal}{\mathcal{G}}
\newcommand{\Hcal}{\mathcal{H}}
\newcommand{\Ical}{\mathcal{I}}
\newcommand{\Jcal}{\mathcal{J}}
\newcommand{\Kcal}{\mathcal{K}}
\newcommand{\Lcal}{\mathcal{L}}
\newcommand{\Mcal}{\mathcal{M}}
\newcommand{\Ncal}{\mathcal{N}}
\newcommand{\Ocal}{\mathcal{O}}
\newcommand{\Pcal}{\mathcal{P}}
\newcommand{\Qcal}{\mathcal{Q}}
\newcommand{\Rcal}{\mathcal{R}}
\newcommand{\Scal}{\mathcal{S}}
\newcommand{\Tcal}{\mathcal{T}}
\newcommand{\Ucal}{\mathcal{U}}
\newcommand{\Vcal}{\mathcal{V}}
\newcommand{\Wcal}{\mathcal{W}}
\newcommand{\Xcal}{\mathcal{X}}
\newcommand{\Ycal}{\mathcal{Y}}
\newcommand{\Zcal}{\mathcal{Z}}

\newcommand{\Abf}{\mathbf{A}}
\newcommand{\Bbf}{\mathbf{B}}
\newcommand{\Cbf}{\mathbf{C}}
\newcommand{\Dbf}{\mathbf{D}}
\newcommand{\Ebf}{\mathbf{E}}
\newcommand{\Fbf}{\mathbf{F}}
\newcommand{\Gbf}{\mathbf{G}}
\newcommand{\Hbf}{\mathbf{H}}
\newcommand{\Ibf}{\mathbf{I}}
\newcommand{\Jbf}{\mathbf{J}}
\newcommand{\Kbf}{\mathbf{K}}
\newcommand{\Lbf}{\mathbf{L}}
\newcommand{\Mbf}{\mathbf{M}}
\newcommand{\Nbf}{\mathbf{N}}
\newcommand{\Obf}{\mathbf{O}}
\newcommand{\Pbf}{\mathbf{P}}
\newcommand{\Qbf}{\mathbf{Q}}
\newcommand{\Rbf}{\mathbf{R}}
\newcommand{\Sbf}{\mathbf{S}}
\newcommand{\Tbf}{\mathbf{T}}
\newcommand{\Ubf}{\mathbf{U}}
\newcommand{\Vbf}{\mathbf{V}}
\newcommand{\Wbf}{\mathbf{W}}
\newcommand{\Xbf}{\mathbf{X}}
\newcommand{\Ybf}{\mathbf{Y}}
\newcommand{\Zbf}{\mathbf{Z}}

\newcommand{\tbf}{\mathbf{t}}
\newcommand{\xbf}{\mathbf{x}}

\newcommand{\Qf}{\mathcal{Q}}

\newcommand{\fa}{\mathfrak{a}}
\newcommand{\fb}{\mathfrak{b}}

%\preceq	\succeq  \precsim

\newcommand{\trv}{conditional random variable}
%\newcommand{\trv}{random field}

\newcommand{\track}{track}
\newcommand{\trackable}{trackable}
\newcommand{\trackability}{trackability}
%\newcommand{\Trackable}{Trackable}
\newcommand{\Trackability}{Trackability}

\newcommand{\sm}{\setminus}						
\newcommand{\ins}{\subseteq} 					
\newcommand{\sni}{\supseteq} 					
\newcommand{\cmpl}{\mathsf{c}}

\newcommand{\ol}{\overline}
\newcommand{\ul}{\underline}

\newcommand{\pr}{\mathrm{pr}} 	
\newcommand{\ev}{\mathrm{ev}} 

\newcommand{\srj}{\twoheadrightarrow}
\newcommand{\inj}{\hookrightarrow}
\newcommand{\dshto}{\dashrightarrow}
\newcommand{\dshot}{\dashlefttarrow}

\renewcommand{\Pr}{\mathbb{P}} 	
\newcommand{\Var}{\mathrm{Var}}
\newcommand{\Bias}{\mathrm{Bias}}
\newcommand{\Noise}{\mathrm{Noise}}
\DeclareMathOperator{\Tr}{Tr}

\newcommand{\Msf}{\mathcal{M}}
\newcommand{\Prob}{\mathcal{P}}

	
\newcommand{\E}{\mathbb{E}}
\DeclareMathOperator{\CE}{\mathrm{CE}}
\DeclareMathOperator{\KL}{\mathrm{KL}}
\DeclareMathOperator{\TC}{\mathrm{TC}}
\newcommand{\I}{\mathbbm{1}}
\newcommand{\Pa}{\mathrm{Pa}} 		
\newcommand{\PaD}{\mathrm{PaD}} 		
\newcommand{\Fa}{\mathrm{Fa}} 		
\newcommand{\Pae}{\mathrm{Pae}} 		
\newcommand{\Pai}{\mathrm{Pai}} 		
\newcommand{\pa}{\mathrm{pa}} 		
\newcommand{\Ch}{\mathrm{Ch}} 	
\newcommand{\Sib}{\mathrm{Sib}} 		
\newcommand{\Adj}{\mathrm{Adj}} 		
\newcommand{\Anc}{\mathrm{Anc}} 		
\newcommand{\AnCl}{\mathrm{AnCl}} 
\newcommand{\Desc}{\mathrm{Desc}} 	
\newcommand{\Dist}{\mathrm{Dist}} 
\newcommand{\Di}{\mathrm{Di}} 
\newcommand{\Dst}{\mathcal{D}} 
\newcommand{\IDC}{\mathrm{ID}}
\newcommand*{\idc}[1]{\breve{#1}}
\newcommand{\eDist}{\mathrm{eDist}} 
\newcommand{\Pred}{\mathrm{Pred}} 
\newcommand{\Sc}{\mathrm{Sc}}
\newcommand{\MBl}{\mathrm{Mb}}
\newcommand{\NonDesc}{\mathrm{NonDesc}}
%\DeclareMathOperator*{\Indep}{\perp\!\!\!\perp}
\DeclareMathOperator*{\Indep}{\perp\mkern-11mu\perp}
\DeclareMathOperator*{\nIndep}{\not\mkern-2mu\perp\mkern-11mu\perp}
%\DeclareMathOperator*{\nIndep}{\not\Indep}
\DeclareMathOperator*{\Perp}{\perp}
\DeclareMathOperator*{\nPerp}{\not\perp}
\DeclareMathOperator*{\given}{|}
\DeclareMathOperator*{\diag}{\mathrm{diag}}
\DeclareMathOperator{\doit}{do}
%\DeclareMathOperator*{\bigdcup}{\dot{\bigcup}}  %error: puts a 9 on top instead of a dot (??)
%\DeclareMathOperator{\dcup}{\dot{\cup}}
%\DeclareMathOperator*{\bigdcup}{\coprod}
%\DeclareMathOperator{\dcup}{\sqcup}
\newcommand{\dcup}{\,\dot{\cup}\,}
\newcommand{\bigdcup}{\mathop{\dot{\bigcup}}}

\DeclareMathOperator*{\argmax}{argmax}
\DeclareMathOperator*{\argmin}{argmin}

\newcommand{\moral}{\mathrm{mor}}	
\newcommand{\marg}{\mathrm{mar}}
%\newcommand{\swig}{\mathrm{o/do}}
\newcommand{\swig}{\mathrm{swig}}
\newcommand{\aug}{\mathrm{aug}}		
\newcommand{\can}{\mathrm{can}}	
\newcommand{\ass}{\mathrm{ass}}		
\newcommand{\augacy}{\mathrm{augacy}}
\newcommand{\acag}{\mathrm{acag}}			
\newcommand{\acy}{\mathrm{acy}}	
\newcommand{\ske}{\mathrm{ske}}	
\newcommand{\cx}{\mathrm{cx}}	
\newcommand{\HEDG}{HEDG}
\newcommand{\Asterisk}{\mathop{\scalebox{1.5}{\raisebox{-0.2ex}{$\ast$}}}}%
\newcommand{\asterisk}{*}

\newcommand{\ReLU}{\mathrm{ReLU}} 
\newcommand{\Lip}{\mathrm{Lip}} 

\newcommand{\lp}{\left ( }
\newcommand{\rp}{\right ) }
\newcommand{\lB}{\left [ }
\newcommand{\rB}{\right ] }
\newcommand{\lC}{\left \{ }
\newcommand{\rC}{\right \} }
\newcommand{\lI}{\left| }
\newcommand{\rI}{\right| }

\newcommand{\st}{\,\middle|\,}

%\newcommand{\tuh}{\rightarrow}
%\newcommand{\hut}{\leftarrow}
%\newcommand{\huh}{\leftrightarrow}
%\newcommand{\tut}{-}
%\newcommand{\tnt}{\diameter}
%\newcommand{\ot}{\leftarrow}
%\newcommand{\oto}{\leftrightarrow}

\newcommand{\arrhead}{{Latex}}
%\newcommand{\arrtail}{{stealth reversed}}  % too visually confusing given al the different edge marks we need
\newcommand{\arrtail}{{}}
\newcommand{\arrstar}{Rays[n=6]}
\newcommand{\arrcirc}{{Circle[open]}}
\newcommand*{\out}[1][]{\mathrel{\tikz [baseline=-0.25ex,\arrcirc-\arrtail, #1] \draw [#1] (0pt,0.5ex) -- (1.3em,0.5ex);}}
\newcommand*{\hut}[1][]{\mathrel{\tikz [baseline=-0.25ex,\arrhead-\arrtail, #1] \draw [#1] (0pt,0.5ex) -- (1.3em,0.5ex);}}
\newcommand*{\tut}[1][]{\mathrel{\tikz [baseline=-0.25ex,\arrtail-\arrtail, #1] \draw [#1] (0pt,0.5ex) -- (1.3em,0.5ex);}}
\newcommand*{\tuh}[1][]{\mathrel{\tikz [baseline=-0.25ex,\arrtail-\arrhead, #1] \draw [#1] (0pt,0.5ex) -- (1.3em,0.5ex);}}
\newcommand*{\tuo}[1][]{\mathrel{\tikz [baseline=-0.25ex,\arrtail-\arrcirc, #1] \draw [#1] (0pt,0.5ex) -- (1.3em,0.5ex);}}
\newcommand*{\huh}[1][]{\mathrel{\tikz [baseline=-0.25ex,\arrhead-\arrhead, #1] \draw [#1] (0pt,0.5ex) -- (1.3em,0.5ex);}}
\newcommand*{\ouo}[1][]{\mathrel{\tikz [baseline=-0.25ex,\arrcirc-\arrcirc, #1] \draw [#1] (0pt,0.5ex) -- (1.3em,0.5ex);}}
\newcommand*{\huo}[1][]{\mathrel{\tikz [baseline=-0.25ex,\arrhead-\arrcirc, #1] \draw [#1] (0pt,0.5ex) -- (1.3em,0.5ex);}}
\newcommand*{\ouh}[1][]{\mathrel{\tikz [baseline=-0.25ex,\arrcirc-\arrhead, #1] \draw [#1] (0pt,0.5ex) -- (1.3em,0.5ex);}}

\newcommand*{\ous}[1][]{\mathrel{\tikz [baseline=-0.25ex,\arrcirc-\arrstar, #1] \draw [#1] (0pt,0.5ex) -- (1.3em,0.5ex);}}
\newcommand*{\hus}[1][]{\mathrel{\tikz [baseline=-0.25ex,\arrhead-\arrstar, #1] \draw [#1] (0pt,0.5ex) -- (1.3em,0.5ex);}}
\newcommand*{\tus}[1][]{\mathrel{\tikz [baseline=-0.25ex,\arrtail-\arrstar, #1] \draw [#1] (0pt,0.5ex) -- (1.3em,0.5ex);}}
\newcommand*{\sut}[1][]{\mathrel{\tikz [baseline=-0.25ex,\arrstar-\arrtail, #1] \draw [#1] (0pt,0.5ex) -- (1.3em,0.5ex);}}
\newcommand*{\suh}[1][]{\mathrel{\tikz [baseline=-0.25ex,\arrstar-\arrhead, #1] \draw [#1] (0pt,0.5ex) -- (1.3em,0.5ex);}}
\newcommand*{\suo}[1][]{\mathrel{\tikz [baseline=-0.25ex,\arrstar-\arrcirc, #1] \draw [#1] (0pt,0.5ex) -- (1.3em,0.5ex);}}
\newcommand*{\sus}[1][]{\mathrel{\tikz [baseline=-0.25ex,\arrstar-\arrstar, #1] \draw [#1] (0pt,0.5ex) -- (1.3em,0.5ex);}}

\newcommand*{\tnt}[1][]{\mathrel{\tikz [baseline=-0.25ex,-, #1] \draw [#1] (0pt,0.5ex) -- node[strike out,draw,-]{} (1.3em,0.5ex);}}
\newcommand*{\toh}[1][]{\mathrel{\tikz [baseline=-0.25ex,-\arrhead, #1] \draw [
decoration={markings, mark=at position 0.4 with {\draw[fill] circle (0.8pt);}},
        postaction={decorate},#1] (0pt,0.5ex) -- (1.3em,0.5ex);}}
\newcommand*{\hot}[1][]{\mathrel{\tikz [baseline=-0.25ex,\arrhead-, #1] \draw [
decoration={markings, mark=at position 0.6 with {\draw[fill] circle (0.8pt);}},
        postaction={decorate},#1] (0pt,0.5ex) -- (1.3em,0.5ex);}}
\newcommand*{\tot}[1][]{\mathrel{\tikz [baseline=-0.25ex,-, #1] \draw [
decoration={markings, mark=at position 0.5 with {\draw[fill] circle (0.8pt);}},
        postaction={decorate},#1] (0pt,0.5ex) -- (1.3em,0.5ex);}}				
\newcommand*{\hoh}[1][]{\mathrel{\tikz [baseline=-0.25ex,\arrhead-\arrhead, #1] \draw [
decoration={markings, mark=at position 0.5 with {\draw[fill] circle (0.8pt);}},
        postaction={decorate},#1] (0pt,0.5ex) -- (1.3em,0.5ex);}}
%\newcommand*{\ouo}[1][]{\mathrel{\tikz [baseline=-0.25ex,-, #1] \draw [decoration={markings, mark=at position 0 with {\draw circle (1pt);}, mark=at position 1 with {\draw circle (1pt);}}, postaction={decorate},#1] (0pt,0.5ex) -- (1.3em,0.5ex);}}

%further arrow tips: Arc Barb, Bar, Bracket, Hooks, Parenthesis, Tee Bar, Circle, Diamond, Rectangle, Square, Rays[n=6]

%\newcommand*{\rars}{\begin{array}{c}\huh\\[-10pt]\tuh\end{array}}
%\newcommand*{\lars}{\begin{array}{c}\huh\\[-10pt]\hut\end{array}}
%\newcommand*{\allars}{\begin{array}{c}\huh\\[-10pt]\tuh\\[-10pt]\hut\\[-10pt]\tut\end{array}}
%\newcommand*{\dars}{\begin{array}{c}\huh\\[-10pt]\tuh\\[-10pt]\hut\end{array}}


\newcommand*{\matb}[1]{\begin{bmatrix}#1\end{bmatrix}}
\newcommand*{\mat}[1]{\begin{pmatrix}#1\end{pmatrix}}


\makeatletter
\newcommand{\xRightarrow}[2][]{\ext@arrow 0359\Rightarrowfill@{#1}{#2}}
\makeatother


\newcommand{\PF}[1]{{\color{orange}{PF: #1}}}
%\newcommand{\PF}[1]{}


\tikzstyle{ndint} = [draw, line width=1pt, color=teal, shape=rectangle, minimum size=20pt,inner sep=0pt, text=black]
\tikzstyle{ndout} = [draw, line width=1pt, shape=circle, color=blue, minimum size=20pt,inner sep=0pt, text=black]
\tikzstyle{ndlat} = [draw, line width=1pt, shape=circle, color=red, minimum size=20pt,inner sep=0pt, text=black]
%\tikzstyle{ndvst} = [draw, ultra thick, shape=swig vsplit, swig vsplit={gap=5pt,line color right=teal, line color left=blue}, inner sep=1.2pt, minimum size=7mm]
%\tikzstyle{ndhst} = [draw, ultra thick, shape=swig hsplit, swig hsplit={gap=5pt,line color lower=teal, line color upper=blue}, inner sep=1.2pt, minimum size=7mm]
\tikzstyle{arint} = [style={->,>=Latex,thick,teal}]
\tikzstyle{arout} = [style={->,>=Latex,thick,blue}]
\tikzstyle{arlout} = [style={->,>=Latex,thick,red}]
\tikzstyle{arlat} = [style={<->,>=Latex,thick,red}]
\tikzstyle{car} = [style={o->,>=Latex,thick,purple}]
\tikzstyle{carc} = [style={o-o,>=Latex,thick,purple}]
\tikzstyle{sarr}=[style={{Rays[n=6]}->,>=Latex,thick}]
\tikzstyle{sars}=[style={{Rays[n=6]}-{Rays[n=6]},thick}]
